\section{Preparing a cell by rank cuts}\label{ap:preparation}

The comparison estimates for the paths of Section~\ref{ap:paths}
require sufficiently high rank of the current determining polynomial. We now
modify the cell until failures of these rank tests have the required
small probabilities. If a failure probability exceeds its prescribed
bound, we obtain a relation at the root and use it to remove one
dimension from a rational value space. The modification must retain the old
integer tuple, since every higher determining polynomial is evaluated
on that tuple and can have arbitrarily large coefficients.

The following proposition is the preparation statement needed for
the forward passage through the cell. Its probability refers to the
stopped tree defined in Section~\ref{ap:path:stopped-tree}. After each modification
we test the new tree; successful tests from an earlier tree are not
assumed to persist.

\begin{proposition}[Prepared system with triangular costs]
\label{ap:major:preparation}
Given a triangular two-box certificate, forecast dimension bounds
and descending numerical budgets in the sense of
Definition~\ref{ap:budgets}.  After polynomially
many rank-cut transactions, there is a certificate with an
arbitrarily small prescribed total proportional target discount
such that every rank-stop probability in the actual current
sampling procedure is at most its prescribed bound.  The
procedure may restart its tests after every transaction.

The dimensions added at block $i$ are polynomial in the dimensions
strictly above $i$.  The increase in its width logarithm is
polynomial in dimension data, $p$, and higher-level numerical
budgets, and is independent of its own initial width logarithm
except for carrying that logarithm additively.  Root slice costs
and sufficient root-side cutoffs have polynomial logarithms in
the full eventual budgets.
\end{proposition}

Throughout this section the degree bound is fixed. Polynomial
dependence on dimensions is permitted. A dimension-controlled factor
means a quantity whose size and inverse size are at most the
exponential of a fixed polynomial in the total relevant dimensions.
The polynomial exponents depend only on the fixed degree bounds.
Warm and cold-pre budgets have the
descending meaning of Definition~\ref{ap:budgets}; a warm bound at
$j$ does not use $Q_i=\log(2/w_i)$ for $i\le j$. All finite
families of requests are forecast before the late perturbation,
range, rank, and side choices.

\subsection{The relation needed for a rank cut}


Return to a triangular system with blocks $b_i$ of weight $i$ and
$C_i$ taking values in rational spaces $W_i$.  Put
\[
                         A_i=[C_i]_{\mathrm{wt}=i}.
\]
Write $\upsilon_i$ for the physical lift-chart width at layer $i$,
the quantity denoted by $v$ in Section~\ref{ap:paths}.
Let $K_i\supset W_i$ be the space obtained by setting the omitted
sheet coordinates at level $i$ to zero.  These spaces and their
coordinate classifications are fixed for the current candidate cell,
independently of the sampled paths.  Their log heights are cold-pre
at their own levels: an omitted integral row has norm at most a warm
factor times $\upsilon_i^{-1}$.

\begin{lemma}[Backward rank transfer]\label{ap:major:backward}
Suppose the actual stopped sampling procedure fails the rank test
at block $h$ with probability larger than the prescribed small
bound at that block.  Then there is an integral row $\lambda$,
nonzero on $W_h$ and of norm at most $R_h$, for which, on
$\R^n\times\bigoplus_{i<h}K_i$,
\begin{equation}\label{ap:major:root-relation}
 \lambda A_h=
 S\left(\frac{u}{\widetilde H},
       \left(\frac{b_i-A_i(u,b_{<i})}{w_i\varepsilon_i}\right)_{i<h}
       \right)+R(u,b_{<h}).
\end{equation}
The log coefficient bound for $S$ and the log common denominator of
$R$ in the raw integer variables are polynomial in the cold budget
at $h$ alone, including dimensions.  Each
$\log(\varepsilon_i^{-1})$ is warm at $i$, and the spatial shrinkages
from $H$ to $\widetilde H$ have polynomial log cost.  The slow and
rational terms have ambient extensions with these respective bounds.
The identity is asserted on the displayed subspaces, not outside them.
\end{lemma}

We first prove the geometric consequence of this relation. This
identifies the two estimates that its proof must provide: the real
term must be controlled in normalized residual coordinates, whereas
the rational term needs a common denominator in integer coordinates.
The proof of Lemma~\ref{ap:major:backward} then occupies the next
three subsections.

\subsection{A rank cut preserving the box certificate}

We now turn the relation into a geometric transaction.  At every
site the strict and enlarged tests count integral tuples satisfying
the triangular residual conditions.  The radii are small enough for
uniqueness, including all the fixed buffers used below.  We retain
the actual tuple throughout; higher polynomials are never estimated
by continuity in an integer slot.

\begin{lemma}[Relative rank-cut transaction]
\label{ap:major:transaction}
Suppose a triangular two-box certificate
\[
                         \E f B^- > a\E B^+
\]
has a relation \eqref{ap:major:root-relation} with $\lambda$ nonzero
on $W_h$.  After a spatial slice, lower residual refinements, and
a prescribed small proportional discount of $a$, there is another
certificate with the following properties.
\begin{enumerate}
\item The new polynomial at $h$ takes values in
$W_h\cap\ker\lambda$.  If $h>1$, at most one full copy of $\R^{d_h}$
is added at weight $h-1$; if $h=1$, there is no incoming copy.
\item At weight $h$, the relative width loss uses only dimension
and gap factors, independently of $R_h$, $Q_h$, and the size of
the rational section used in the cut.  At weights above $h$, no
width loss is necessary.  Lower relative losses are polynomial-log
in their warm data and the cold data at $h$.
\item Every old higher polynomial is transformed by an exact
weighted polynomial substitution on its actual integer arguments.
All weights remain valid.  Rational value spaces are preserved,
apart from the stated intersection and full incoming copy.
\item Spatial slice costs and sufficient root-side cutoffs have
polynomial logarithms in the full budgets.  There is no loss
proportional to the inverse of the actual mass of the old test.
\end{enumerate}
\end{lemma}

\begin{proof}
Choose a rational $e\in W_h$ with $\lambda e=1$, and set
\[
                   C_h^{\mathrm{new}}=(1-e\lambda)C_h.
\]
All subsequent changes of lower variables are composed into this
formula.  Its values belong to $W_h\cap\ker\lambda$ exactly.
We first construct, on suitable lower and spatial cells, the identity
\begin{equation}\label{ap:major:cut-identity}
 C_h^{\mathrm{old}}=
 C_h^{\mathrm{new}}+C_{\mathrm{in}}+c+I+E,
                     \qquad \|E\|_\infty\ll\gamma w_h.
\end{equation}
Here $I$ is integral-valued on the reindexed variables and has
weight at most $h$; $c$ is constant; and, when $h>1$,
$C_{\mathrm{in}}$ has weight at most $h-1$ and depends only on
arguments strictly below $h-1$.

Prescribe the spatial and lower integer variables modulo a common
integer clearing the denominator of $eR$ in
\eqref{ap:major:root-relation}.  Its log size is polynomial in the
cold budget at $h$ and the section bound.  On reindexing a residue
class, every positive-degree coefficient of this rational polynomial
becomes integral; its constant is absorbed into $c$.  Arbitrarily
large numerators do not affect this assertion.

Subdivide the lower residuals by warm relative factors.  Each omitted
sheet row annihilates $C_i$ and is integral on $b_i$.  Its norm is
at most a warm factor times $\upsilon_i^{-1}$, so it varies by less than
$1$ on a sufficiently small enlarged lower cell.  It is therefore
constant there.  Thus each nonempty enlarged lower cell puts $b_i$
on a fixed affine translate of $K_i$.

First take the weight-$h$ parts of both terms in
\eqref{ap:major:root-relation}. The left side is homogeneous of
weight $h$, and each deviation $b_i-A_i$ is homogeneous of weight
$i$, so this preserves the identity and does not enlarge either
the coefficient bound or the common denominator. Thus we may
take $S$ and $R$ homogeneous of weight $h$ in their respective
weighted coordinates.

Let $u_*$ be the spatial cell center, and let $r_i$ be the center
of the chosen lower box in the residual coordinate $b_i-C_i$.
Define the completed slow polynomial by
\[
 \widehat S(u,b_{<h})=
 S\left(\frac{u-u_*}{\widetilde H},
   \left(\frac{b_i-C_i(u,b_{<i})-r_i}
                   {w_i\varepsilon_i}\right)_{i<h}\right).
\]
Since $C_i$ has top weighted part $A_i$, these replacements and
constant translations preserve the top weighted symbol. In the
joint spatial and residual coordinates $\widehat S$ has controlled
coefficients. Refining those coordinates
makes its oscillation, even after multiplication by $e$, much less
than $\gamma w_h$.  The additional relative precision uses only
dimension factors, $\gamma^{-1}$, $w_h^{-1}$, $\|e\|$, and the
coefficient bound in \eqref{ap:major:root-relation}.  For every lower
level these are permitted higher-level costs.  No derivative bound
for the original $C_i$ is needed: the residual coordinates themselves
are being localized.

The difference between $\lambda C_h$ and this completed slow term
plus $R$ has vanishing weight-$h$ part on
$\R^n\times\bigoplus_{i<h}K_i$.  Off that subspace it need not have
smaller weight.  Compose the difference with rational affine
projections onto the affine spaces of the chosen cell.  This gives
a formal polynomial of weight at most $h-1$ agreeing with the
difference on the whole cell.  Multiplying by $e$ produces a
lower-weight remainder, while the slow term contributes a constant
and the indicated small error and the rational term contributes
$I$ and a constant.

Implement a lower residue restriction by
$b_i=M b'_i+b_{i0}$ and use $C'_i=C_i/M$ after the corresponding
substitutions; put the constant offset in the residual center.
This preserves $W_i$.  If $h>1$, the lower-weight remainder has
only a constant-coefficient linear dependence on the variables of
weight $h-1$.  Reduce that coefficient matrix modulo integers,
putting the removed integral term into $I$.  Its remaining entries
are bounded.  Replace $b'_{h-1}$ there by $C'_{h-1}$ plus the chosen
residual center.  A further lower refinement to width
$\ll\gamma w_h/(1+d_{h-1})$ makes the error suitably small.  What
remains is $C_{\mathrm{in}}$, depending only on lower arguments.
It need not take values in the line spanned by $e$, which is why
we allow a full $d_h$-dimensional incoming copy.  When $h=1$, the
remainder is constant instead.  This proves
\eqref{ap:major:cut-identity}.

Choose a narrow residual box modulo $1$ for a new integer block
$b_{\mathrm{in}}$ around $C_{\mathrm{in}}$, with radius much less
than $\gamma w_h$.  Introduce $b_h^{\mathrm{new}}$ by the exact identity
\begin{equation}\label{ap:major:tuple-transport}
             b_h^{\mathrm{old}}
                 =b_h^{\mathrm{new}}+b_{\mathrm{in}}+I,
\end{equation}
with a possible constant integer adjustment.  Apply this exact
substitution in every $C_i$ for $i>h$.  Since $I$ has weight at
most $h$ and $b_{\mathrm{in}}$ has weight $h-1$, their weighted
degree bounds are preserved.  No continuity estimate in an old
integer argument occurs, and all higher widths can be kept.

We give the selection argument, including its uniform width control.
At lower residual boxes use continuously shifted smaller boxes,
allowing centers in a small neighborhood of the old strict box but
well inside its enlargement.  Every old strict point receives the
full smaller strict coverage volume, and every new enlarged point
lies in the old enlarged box.  The same construction works after
the uniform residue scaling.

For the incoming block choose its center $l_{\mathrm{in}}$ uniformly
in a unit cube.  For any old strict tuple, the total center measure
of admissible incoming strict choices is exactly their box volume,
including integer indexing across wrap.  Let the new $h$-center
range within radius $(1+\gamma/2)w_h$ about
$l_h+c-l_{\mathrm{in}}$.  Write
$r_{\mathrm{in}}=b_{\mathrm{in}}-C_{\mathrm{in}}$ and
$r_h^{\mathrm{new}}=b_h^{\mathrm{new}}-C_h^{\mathrm{new}}$.
Equations~\eqref{ap:major:cut-identity} and
\eqref{ap:major:tuple-transport} give
\[
       b_h^{\mathrm{old}}-C_h^{\mathrm{old}}
             =r_h^{\mathrm{new}}+r_{\mathrm{in}}-c-E.
\]
Choose the two new radii and $\|E\|_\infty$ well below
$\gamma w_h/8$.  An old strict tuple then receives full strict
coverage in the new $h$-center, and every transported enlarged tuple
lies in the old enlarged $h$-box.  The new $h$ radius can be a fixed
dimension/gap fraction of $w_h$, independently of $e$ or $R_h$.
Within weight $h-1$ take a common decreased width for the existing
and incoming coordinates.  This may make the incoming box
arbitrarily smaller than $\gamma w_h$, but does not require the
error in \eqref{ap:major:cut-identity} to be small at that new
incoming scale: the error is used only in the old $h$-gap.

Let $c_-$ be the product of the full strict coverage volumes and
$c_+$ the corresponding enlarged upper coverage constant, including
the residue and spatial averaging.  For every site, counting its
exact tuples gives
\[
 \int B^-_{\mathrm{new}}\ge c_- B^-_{\mathrm{old}},
       \qquad
 \int B^+_{\mathrm{new}}\le c_+ B^+_{\mathrm{old}}.
\]
The first inequality persists after multiplication by $f\ge0$.
The ratio satisfies
\[
                  c_+/c_-\le
                (1+O(\gamma))^{O(\sum_i d_i+d_h)}.
\]
Unchanged coordinates have ratio $1$ in this coverage calculation.
Consequently, with $a'=a c_-/c_+$, integration of
$\E fB^-_{\mathrm{new}}-a'\E B^+_{\mathrm{new}}$ is strictly
positive.  One cell gives the desired certificate.  There is no
division by the actual mass of either old test.  Sum over residue
classes and average spatial slices with their actual box masses;
equal or nearly equal subdivisions have controlled relative sides
once the parent cutoff is imposed.  All constants, affine
projections, and incoming polynomials may depend on the selected
lower cell.  They are measurable, for example by using fixed
projections for each discrete sheet value.

Finally shift centers and integer indices into their bounded
representative ranges.  This only changes constant integer
adjustments in \eqref{ap:major:tuple-transport}, not a value space.
Every asserted cost follows from the refinements above.  Large
section norms are charged only to lower/spatial precision; the own
layer width uses only its gap and dimensions.  This proves all
parts of the lemma.
\end{proof}

\subsection{Two recovery estimates}

Consider one invocation of the sampler, at level $j$, satisfying
the parent rank and side hypotheses of Lemma~\ref{ap:path:ensemble}
at all the requested accuracies, with ordinary
polynomial $C:U\to W\subset E=\R^d$, and write $A=C_{[j]}$.  The sides
of $U$ are $H_i$, and the common lift-chart width is $v$.  Write
$(w,z)$ for the unimodular integer coordinates on the lift integer
$\beta$, with $z$ the projected sheet coordinates.  The corresponding
projected lattice basis vectors are denoted by $v_i$, and
$k_i=\max(1,K_i v)$, with the harmless rounded convention of the
sampler.  No relation between this symbol $w$ and a block width is
intended.

Remove every omitted sheet direction and, if needed, some additional
inactive directions.  Removing a direction always means removing
\emph{all} its parameter columns, including its additional linear
column.  Let $K\supset W$ be the subspace on which the removed sheet
coordinates vanish.  The leading substitution after these columns
have been set to zero is
\[
 \Phi(t)=\bigl(\psi_{[1]}(t),\beta_{[j]}(t)\bigr)
       \in\R^n\times K.
\]
We use the full primitive lattice $K\cap\Z^d$, not an arbitrary
rationally scaled basis of $K$.

The real recovery estimate must be independent of the unrestricted
coefficients of $A$. We therefore compare a lift variable with its
predicted value $A(u)$ before measuring coefficients. This is
different from choosing an integer basis for $K$, whose real
conditioning may be poor.

\begin{lemma}[Normalized real recovery]\label{ap:major:real-recovery}
Let $a_j$ be the common normalized spatial scale of the kernel blocks.
On the space of weighted homogeneous polynomials of any prescribed
bounded degree, substitution by $\Phi$ has a bounded inverse on its
image when coefficients on $\R^n\times K$ are measured in
\begin{equation}\label{ap:major:units}
       \left(\frac{u_i}{a_jH_i}\right)_i,
       \qquad \frac{b-A(u)}{v},
\end{equation}
and coefficients on the parameter box are measured in $t_\ell/T_\ell$.
Both the forward operator norm and the inverse norm are bounded by
dimension-controlled factors.  The assertion holds for coefficient
tensors in an additional real normed coefficient space.

If all normalized parameter variables are restricted to radius $\zeta$,
replace $a_j$ by $a_j\zeta$ and $v$ by $v\zeta^j$ in
\eqref{ap:major:units}.  The same conclusion holds.  The constants do
not depend on the widths or on the coefficients of $C$.
\end{lemma}

\begin{proof}
The homogeneous identity supplied by the lift is
\[
                   b-A(u)=-y_{[j]}.
\]
Consequently the unbounded coefficients of $A$ disappear after the
change of coordinates \eqref{ap:major:units}.  In these coordinates
the ideal map has spatial kernel blocks and lift principal blocks
supported on disjoint parameter columns.  Fix all but one variable
in a suitable dedicated product for each retained lift direction,
at points bounded away from zero in their normalized ranges.  Together
with an invertible spatial kernel block, the remaining variables
then give a full-dimensional affine image whose sides are bounded
below by dimension-controlled constants.

For an inactive axis, the product coefficient is $1$, and the product
of the sides is comparable to $k_i$ within dimension-controlled
factors.  The extra linear column can be set to zero; this includes
the case $j=1$.  Axes whose rounded scale could obstruct this
comparison have already been omitted by the minimum-side convention.
The split-basis conditioning for the remaining lift directions is
dimension-controlled by the lattice chart construction.  Thus a fixed
box in the real coordinates \eqref{ap:major:units} is covered with a
dimension-controlled inverse parametrization.

At fixed degree, interpolation on this box bounds every coefficient
by a dimension-controlled multiple of the supremum of the polynomial
on the box.  Conversely, the normalized principal substitution has
dimension-controlled coefficients, so it bounds the coefficients of
every pullback.  There are only polynomially many coefficients at
fixed degree.  These two observations give the asserted norms on the
finite-dimensional coefficient spaces.  They also apply with a
vector-valued coefficient, by the same interpolation formulas.

The perturbations are chosen small in these normalized units.
The resulting operator on the coefficient space is therefore a small
perturbation of an injective operator with the displayed inverse
bound.  Choosing the late perturbation sufficiently small preserves
the bounds.  Deleting entire inactive column families does not alter
the retained principal blocks.  Finally, the leading spatial map has
degree $1$ and the leading lift map has degree $j$, so rescaling all
parameters by $\zeta$ gives exactly the two stated powers of $\zeta$.
\end{proof}

The arithmetic recovery estimate uses the extra pure-free columns.
Their structural number is chosen before any accuracy budget.  On a
separate pure-free unit axis, the spatial degree-$1$ coefficients and
the degree-$j$ deck and active-sheet coefficients form one output
vector.  The coefficient families used for different such axes are
disjoint.  Lemma~\ref{ap:path:ensemble} makes these
vectors jointly uniform modulo any one prescribed modulus, to an
arbitrarily small preassigned error.  In the abstract Haar-cover
comparison law, deck residues are uniform even conditional on the
sheet data.  For the actual sampler we use the resulting joint
finite-residue comparison, not conditional uniformity given exact
real sheet coefficients.  On a moderate axis the added
top-degree integer coefficient range has size comparable to
$k_i\sigma/L^j\ge A^j\sigma$; hence the late scale gap gives the same
conclusion for these coefficients.  Spatial free columns have broad
integer supports as well.  Only separate modulus tests, not a test
at the product of all requested moduli, are used below.

\begin{lemma}[Prime-power polynomial test]\label{ap:major:padic}
For each fixed $r$ there are constants $c_r,C_r>0$ such that, for every
prime $p$, integer $e\ge1$, and polynomial $P$ over $\Z/p^e\Z$ of
total degree at most $r$ with at least one unit coefficient,
\[
       \Pr\bigl(P(X)=0\pmod {p^e}\bigr)
               \le C_r(p^e)^{-c_r},
\]
where $X$ has independent uniform coordinates.  The constants are
independent of the number of variables.
\end{lemma}

\begin{proof}
First consider one variable, allowing a coefficient of valuation at
most $(1-\eta)e$.  Put $B=\binom{r+1}{2}$ and, for $r\ge1$, choose
$\ell=\lceil\eta e/(2B)\rceil$.  If the zero set occupied more than
$r$ classes modulo $p^\ell$, choose $r+1$ zeros in different such
classes.  Their Vandermonde determinant has valuation at most
$(\ell-1)B<\eta e/2$.  Multiplying the evaluation vector by the
adjugate shows that every coefficient has valuation at least
$e-(\ell-1)B>(1-\eta)e$, a contradiction.  Thus the zero set has
measure at most $r p^{-\ell}$.  Bounded depths and degree zero are
absorbed into constants.

The same calculation proves the following induction statement: for
an $m$-variable polynomial with a specified coefficient of valuation
$v<e$, its zero probability is at most
\[
                   mr\,p^{-(e-v)/(2B\,2^{m-1})}.
\]
For the induction step, use the coefficient polynomial in the first
$m-1$ variables containing the specified coefficient and test it
to depth $h=\lceil(e+v)/2\rceil$.  The induction hypothesis bounds
its exceptional probability using $h-v\ge(e-v)/2$.  Outside that
set its value has valuation at most $h-1$, and the univariate
calculation has gap $e-h+1\ge(e-v)/2$.  The constants add to at
most $mr$.  A constant coefficient polynomial has empty exceptional
set and satisfies the same estimate.

For many variables, start with a monomial having a unit coefficient.
If a strictly componentwise dominating monomial has coefficient
valuation below the midpoint of the present valuation and $e$, move
to that monomial.  Each move strictly increases total degree, so
there are at most $r$ moves.  At termination the chosen coefficient
has valuation at most $(1-2^{-r})e$, and every strictly dominating
coefficient has larger valuation, separated from it by a fixed
degree-dependent fraction of the remaining depth.

Fix all variables not occurring in the chosen monomial.  In the
coefficient of that monomial, every new contribution comes from a
strictly dominating monomial, and hence cannot cancel the chosen
coefficient.  Its valuation is unchanged.  There remain at most
$r$ variables.  If the chosen monomial is constant, all nonconstant
contributions have larger valuation and the same observation already
settles the claim.  Otherwise use the displayed induction statement
with $m\le r$ and $e-v\ge2^{-r}e$.  For example, the deliberately
weak choices $C_r=r^2$ and $c_r=1/(2^{2r}B)$ suffice for $r\ge1$.
\end{proof}

\begin{lemma}[Simultaneous denominator recovery]
\label{ap:major:denominator}
Prescribe a cold-pre upper bound $H$ for the least common denominator
of the coefficients of a rational polynomial of bounded degree, and
an exceptional probability $\delta>0$.  The late sampler parameters
can be chosen so that, outside an event of probability at most
$\delta$, there is an integer $D_0$ with
\[
          \log D_0\le
          \operatorname{poly}\bigl(d+n+\log(2/\delta)\bigr)
\]
having the following property simultaneously for every allowed
subspace $K$.  If $F$ is weighted homogeneous on $\R^n\times K$, its
coefficients in the primitive lattice coordinates have least common
denominator at most $H$, and $d'$ is a positive integer for which
$d'F\circ\Phi$ has integral coefficients,
then $d'D_0F$ has integral coefficients in those coordinates.
There is no upper-bound requirement on $d'$.  The same assertion
holds for polynomially many components and formal coefficient
tensors with a common initial denominator.  The bound on $D_0$ is
independent of $H$.
\end{lemma}

\begin{proof}
First set the inactive parameter columns to zero.  Let $N$ be the
number of scalar monomials in the remaining output coordinates at
the prescribed degree bound; $N$ is polynomial in their dimension.
For $q=p^e$ above a degree-dependent constant, choose sufficiently
many disjoint pure-free axes that Lemma~\ref{ap:major:padic} and a
union bound over at most $q^N$ coefficient vectors give probability
at most $q^{-10}$ that \emph{some} primitive polynomial vanishes at
all their output vectors.  Indeed, if there are $M$ axes, the bound
is
\[
                  q^N(C_rq^{-c_r})^M,
\]
and $M$ can be a fixed polynomial in the structural dimensions.
An identically zero substituted polynomial vanishes at each unit
axis, so the same bound applies to bad polynomial substitution.
The union over all polynomials is what permits $F$ to be chosen
after the sampled path.

For each prime let $b_p$ be the greatest tested exceptional depth,
also including a fixed baseline for prime powers below the preceding
constant.  Under the uniform residue law, CRT makes bad events at
distinct primes independent.  Thus simultaneous bad powers
$q_1,\ldots,q_t$ at distinct primes have probability at most
$(q_1\cdots q_t)^{-10}$, apart from the fixed baseline.
There is a finite-witness bound uniform in the tested range.  If
$\prod_p p^{b_p}>R$, either one factor exceeds $R$, or a subproduct
lies in $(R,R^2]$.  Summing respectively over prime powers and over
integers gives an $O(R^{-9})$ upper bound, with degree-dependent
constants.  Taking $R$ polynomial in $2/\delta$ bounds this
exceptional product outside probability at most $\delta/2$.

For the actual sampler, use its joint finite-residue law separately
at every required prime power and every product witness in
$(R,R^2]$.  These form a finite family.  The logarithms of their
number and largest modulus are bounded by the pre-bound on $H$ and
$\log(2/\delta)$.  Choose the individual approximation errors so
that their sum is below $\delta/2$.  Thus the input precision may
depend on $H$, but the bound on the exceptional product does not.
This argument uses no approximation at the product of all tested
primes.

On a good path, the following integral consequence holds on the
pure-free outputs.  If substitution makes an integral polynomial
divisible by $p^a$, then all its coefficients are divisible by
$p^{a-b_p}$.  Otherwise divide by its least coefficient valuation;
the resulting primitive polynomial would give an exceptional depth
larger than $b_p$.

Restore inactive outputs by increasing total degree in those
coordinates.  After setting unused nonfree columns and the extra
linear columns to zero, each retained inactive coordinate is a sum
of disjoint dedicated degree-$j$ products, each with coefficient
$1$.  Suppose the least inactive degree not yet removed is $a_{\rm in}$,
and write its terms as
$\sum_{|\alpha|=a_{\rm in}}z^\alpha P_\alpha$ in the other outputs.
Choose one dedicated product for each relevant coordinate and
extract the product with exponent $\alpha_i$ in its parameters.
Its degree in these parameters is $ja_{\rm in}$, and its coefficient is
exactly $P_\alpha$ evaluated on the pure-free outputs.  Higher
inactive degrees cannot contribute.  Dependence of other outputs
on these dedicated parameters adds positive parameter degree and
also cannot contribute at degree $ja_{\rm in}$.  No factorial division is
involved.

Consequently, if the whole substituted integral polynomial is
divisible by $p^a$, this extraction recovers the coefficients at the
least inactive degree with loss at most $b_p$.  Subtract the
corresponding integral coefficient terms.  Their substitutions
retain that divisibility, so the residual is divisible to the
reduced depth.  Continue through the bounded number of possible
inactive degrees.  The total loss is at most a degree-dependent
multiple of $b_p$.  The same argument works for any subset of
removed inactive directions with the same pure-free goodness; it
does not require a new probability estimate for every subspace.

Finally let $d$ be the least common denominator of $F$ and put
$P=dF$.  For a prime dividing $d$, set
$e=\max(0,v_p(d)-v_p(d'))$.  The hypothesis implies that
$P\circ\Phi$ is divisible by $p^e$, and $e\le\log_p H$.
Since $P$ has a coefficient prime to $p$, the preceding divisibility
conclusion forces $e\le C b_p$ for a fixed-degree constant $C$.
Thus $d$ divides $d'(\prod_p p^{b_p})^C$.  Absorb this fixed power
into $D_0$.  Apply the same argument to scalar tensor entries; their
simultaneous polynomial test and common denominator give the claimed
joint conclusion.
\end{proof}

\subsection{The sharpened major statement}

The next lemma permits preliminary restrictions to entire short
inactive directions.  Their lengths, and the height of the subspace
they determine, are part of the cold-pre input.  The important output
is sharper than that input budget.

\begin{lemma}[Relative major refinement]\label{ap:major:refinement}
Let $D(u,b)$ be a fixed weighted homogeneous scalar polynomial of
bounded degree on $\R^n\times K$, with no coefficient bound.  Suppose
that on at least a prescribed cold-pre fraction of paths its full
leading substitution has a slow-plus-rational decomposition: the
slow coefficient norm on the normalized parameter box of radius
$\zeta$ is at most $B$, and the rational term has a common
denominator at most $\rho$.  These input bounds, the inverse fraction,
and $\zeta^{-1}$ have cold-pre logarithms.  Assume the paths have
Lemma~\ref{ap:major:denominator}'s goodness, with their respective
path-dependent multipliers bounded by $D_0$,
through all denominator bounds forecast for the argument.

Then on $\R^n\times K$,
\begin{equation}\label{ap:major:sharp-form}
 D(u,b)=
 S\left(\left(\frac{u_i}{H_i a_j\zeta}\right)_i,
               \frac{b-A(u)}{v\zeta^j}\right)+R(u,b),
\end{equation}
where the coefficient norm of $S$ is at most
$B\exp(\operatorname{poly}(d+n))$, and the logarithm of a common
denominator of $R$ in primitive lattice coordinates is bounded by
\[
       \operatorname{poly}\bigl(d+n+\log(2+D_0+\rho)\bigr).
\]
The assertion is simultaneous for polynomially many coefficient
tensors.  Their real norm may contain a specified pre-scaling, and
their integral coordinates may have specified pre-height.  These
heights do not enter the sharpened real bound.

The two terms can be extended separately to ambient polynomials:
orthogonally in deviation coordinates for $S$, and by a primitive
lattice extension for $R$.  Their identity is asserted only on the
specified subspace.
\end{lemma}

\begin{proof}
We first recover the two terms on a subspace $K'$ for which the
remaining parameter sides are long enough to separate small real
coefficients from rational ones. We then restore the removed
directions by extracting their coefficient-one products. The
restoration is an induction on degree, so its depth is independent
of the number of directions.

Apply Proposition~\ref{lifting:candidate-major} at cold-pre cost to
a fixed ambient extension of $D$, after freezing the already
excluded short coordinates.  The proposition allows unrestricted
coefficients.  Its proof applies to the present sampler: the local
correlation is obtained from a buffered chart-supported phase,
fully slow and integral lift identities from
Lemma~\ref{ap:path:ensemble}, and the cold detection conclusion of
Theorem~\ref{ap:path:detection}.  It
does not use a fine mixed-center spatial boundary estimate.  We
obtain a decomposition with cold-pre bounds on a rational space
$K_0\supset W$ of cold-pre height.

Only after this output, increase the pure-long threshold, deleting
every further inactive direction with a dedicated side below it.
Delete its linear column too; its length is bounded by a fixed
power of a dedicated side, so all these deleted columns are still
pre-bounded.  Denote the resulting subspace of $K$ by $K'$.  Then
$K'\subset K_0$.  Indeed an integral annihilating row of $K_0$
vanishes on $W$.  Its value on a retained projected lattice vector
$v_i$ is an integer, since $v_i$ is the projection of an integer
vector.  Its absolute value is less than $1$ once $K_i$ exceeds
the prescribed height bound, because $|v_i|\ll 1/K_i$.  It is
therefore zero.  This proves the inclusion.

Restrict the candidate decomposition to the remaining pure-long
parameters and compare it with the given decomposition.  In every
positive-degree coefficient their rational difference is the
negative of a slow difference.  The denominator is cold-pre
bounded, and all the remaining sides can exceed the corresponding
separation scale.  Lemma~\ref{prelim:separation}, with zero target
subspace, makes both differences zero.  Hence the candidate slow
and rational terms have, separately, the original bounds $B$ and
$\rho$ on these evaluations.  Apply
Lemma~\ref{ap:major:real-recovery} to the first and
Lemma~\ref{ap:major:denominator} to the second.  This gives the
sharpened conclusions on $K'$.  In degree zero, simply take a split
of the fixed constant from one productive path; no separation of
constants is asserted.

Extend the two sharpened terms separately to $K$.  Orthogonal
projection of a deviation onto $K'$ does not enlarge its real norm.
For the rational term, $K'\cap\Z^d$ is primitive, so a basis of
this lattice extends to an integral basis of the ambient lattice.
Extending the polynomial independently of the complementary
coordinates preserves its denominator.  The resulting numerator
sizes may be arbitrarily large and are not bounded or used in a
slow estimate.  The slow extension retains its real coefficient
bound on the full retained path substitutions, by the forward bound
in Lemma~\ref{ap:major:real-recovery}.  The rational extension needs
only its denominator bound there: integral substitution preserves
it regardless of numerator size.  No real bound on that evaluation
is claimed or needed.  Subtract the two terms.  The
residual vanishes on $K'$.

We finish by induction on the polynomial degree.  Expand the
residual in the extra sheet coordinates for $K/K'$ and take its
least remaining total degree $a\ge1$:
\[
                     \sum_{|\alpha|=a} z^\alpha D_\alpha(u,b'),
                     \qquad b'\in K'.
\]
At this least degree the coefficient tensors are intrinsic; changing
the linear section changes only higher extra-sheet degrees.  Set
the new short linear columns to zero and extract separate dedicated
products of total parameter degree $ja$, as in the denominator
proof.  Each normalized product has size comparable to $k_i\zeta^j$.
Thus the evaluated $D_\alpha$ have slow norm at most the current
norm times
\begin{equation}\label{ap:major:extraction-loss}
        \exp(\operatorname{poly}(d+n))
                      \prod_i(k_i\zeta^j)^{-\alpha_i},
\end{equation}
and rational denominator at the current sharpened bound.  The
coefficient-$1$ extraction introduces no rational division.  Other
outputs' dependence on these parameters contributes only above
degree $ja$.

Each $D_\alpha$ has smaller weighted degree, so apply the present
lemma inductively, simultaneously to its tensor entries, starting
with $K'$.  The newly frozen sides and subspace description are
pre-bounded before this invocation.  Extend the two answers
separately, multiply them by $z^\alpha$, and subtract.  Since
$z_i(A(u))=0$, multiplication by the normalized powers of $z_i$
cancels precisely the losses in \eqref{ap:major:extraction-loss}.
Since the $z_i$ are integral lattice coordinates, the rational
denominator changes only by the common-denominator products from
boundedly many stages.  Continue to the next extra-sheet degree.
The recursion depth is bounded by the fixed polynomial degree, not
by the dimension.  The number of simultaneous components is
polynomial in dimension.  This proves the stated bounds and the
tensor assertion.
\end{proof}

\subsection{Transferring the failed relation through earlier layers}

\begin{proof}[Proof of Lemma~\ref{ap:major:backward}]
There are at most $(1+2R_h)^{d_h}$ possible failed rows.  Pigeonhole
one of them.  Its failure fraction still has inverse logarithm
polynomial in the cold budget at $h$.  Prune to prefixes whose
conditional successful continuation fraction is at least a fixed
power of this fraction; ordinary conditional averaging gives such
prefixes at every level.  All earlier rank tests on them have passed.
For this calculation discard modular-bad lower paths, with an
exceptional product bound $D_0$ whose logarithm is polynomial in the
same budget at $h$.  The exceptional probabilities are small enough
not to remove most of the pruned branch.  The tested denominator
ranges may use the larger cold-pre budgets of those lower layers.
Lemma~\ref{ap:major:denominator} is simultaneous on these ranges, so
its output $D_0$ does not acquire those larger costs.  At each node,
fix one center for which the required conditional fraction remains.
The center does not affect any leading map.

Lift the failed relation backwards, one layer $j<h$ at a time.
The later variables $b_i$, $j<i<h$, are now formal variables on
$K_i$.  Lower histories are fixed at this node.  The induction
hypothesis supplies local slow/rational splittings on the productive
children, with slow coordinates
\[
 \frac{t}{T}\quad\hbox{and}\quad
 \frac{b_i-A_i(\psi_{[1]},\beta_{[j]},b_{>j,<i})}
      {w_i\varepsilon_i},\qquad j<i<h,
\]
allowing a common warm shrink $\zeta$ of the normalized $t$ box.
The initial such splitting is precisely the failed rank certificate.

At each stage we keep two expressions for the same residual: its
raw polynomial in the later integer blocks, and its slow-plus-rational
splitting on productive children. A degree tuple is removed from
both expressions before the next tuple is considered. The slow
coefficient will be multiplied by deviation monomials; the rational
coefficient will be multiplied by raw integer monomials. Before
subtraction, the newly recovered slow coefficient must agree exactly
with the old local slow coefficient, and likewise for the rational
coefficient. A change between those monomials can involve
unrestricted coefficients of the $A_i$.

Order the tuples of degrees in the later blocks lexicographically,
with the highest layer first, in decreasing order.  At a current
tuple $\alpha=(a_i)_{i>j}$, extract the coefficient tensor of the
raw polynomial residual.  Higher tuples have already been removed
from the raw polynomial and from the local splitting.  The local
slow tensor has norm bounded by its prescribed coefficient bound
times $\prod_{i>j}(w_i\varepsilon_i)^{-a_i}$; the local rational
tensor has the accumulated common denominator in primitive
coordinates.  Apply Lemma~\ref{ap:major:refinement} simultaneously
to these coefficients.  Extend its slow tensor orthogonally and its
rational tensor primitively.  Multiply the slow tensor by the
corresponding deviation monomials $(b_i-A_i)^\alpha$, and the rational
tensor by the raw monomials $b_i^\alpha$, and subtract.

We verify the exactness needed for this subtraction.  In positive
remaining degree in $t$, the old local tensor split and the newly
constructed split sum to the same evaluated raw tensor.  Their
slow bounds, expressed on the full normalized $t$ box, use only
dimensions, $Q_{>j}$, and already fixed upper-layer data.  This is
because the forward operator in Lemma~\ref{ap:major:real-recovery}
has a dimension-controlled bound: evaluating a deviation divided
by $\upsilon_j$ does not introduce an inverse power of $\upsilon_j$.  Their
rational denominators have the same allowed dependence.  The
primitive bases for later $K_i$ are also part of the higher-level
data.  The minimum-side convention at $j$ can therefore ensure
separation using only warm information at $j$.  Applying
Lemma~\ref{prelim:separation} coefficientwise shows that the two
slow tensors agree exactly and the two rational tensors agree
exactly.  No assertion of approximate equality is used before
multiplication by a possibly very large coefficient of a later
$A_i$.

In remaining degree zero there is no $u,b_j$ argument.  Pigeonhole
the exact pair of constant tensors instead.  To see that only
pre-many choices occur, compare two such pairs: their rational
difference has bounded denominator, and its size is bounded by
the slow norms, in the fixed later lattice coordinates.  A bounded
region of a bounded-denominator lattice is finite with controlled
log cardinality.  The common center of this region can be arbitrarily
large, which does not affect that count.  The extra fraction loss
is a pre cost for the local application, not a change to the
already imposed modular-goodness multiplier.

These arguments justify subtraction from the local splitting as
well as from the raw polynomial.  They prevent unbounded lower
coefficients arising from the changes between raw and deviation
monomials.  Primitive extensions of homogeneous tensors remain
homogeneous.  Orthogonal extension is compatible with every later
$A_i$, since those values lie in $W_i\subset K_i$.  Continue through
the degree tuples.  There are only boundedly many tuples, because
the number of layers and all degrees are fixed; entries of a tensor
are processed simultaneously rather than one monomial at a time.

Each lift inserts the new deviation scale $\upsilon_j\zeta^j$ and spatial
scale $a_j\zeta$.  Thus its new $\varepsilon_j$ has warm inverse
logarithm at $j$.  The sharpened coefficient and denominator bounds
are changed only by polynomial dimension factors and the bounds
already charged at $h$.  Fixed-depth composition gives
\eqref{ap:major:root-relation}.  When $h=1$, there is no backwards
step and the rank certificate is already this relation.

For completeness, successive local selections do not create a
probability circularity.  At the beginning use the pruned branch
with the selected row and modular goodness.  At any of its nodes
one may subsequently retain a smaller pre fraction to choose a
constant split or perform a major calculation.  The conclusion at
that node is existence of a split with the uniform sharpened bounds.
For the preceding ancestor it is this existence statement that is
used; the smaller internal fraction need not be propagated as an
output coefficient or denominator cost.  All the finitely nested
denominator ranges were forecast before the late sampler choices.
The argument only requires positive conditional fractions at the
nodes where the corresponding major theorem is applied.
\end{proof}

\subsection{Termination and the order of the budgets}

\begin{proof}[Proof of Proposition~\ref{ap:major:preparation}]
If a stop probability is too large, apply
Lemma~\ref{ap:major:backward} at that level and then
Lemma~\ref{ap:major:transaction}.  Retest the resulting actual
system.  Earlier tests are not assumed to remain valid after a
restriction.

To count transactions, let $c_h$ be the number of cuts at level
$h$.  Each such cut removes one dimension of $W_h$.  New dimensions
there arise only from incoming copies at higher levels.  Thus
$c_h$ is at most the initial dimension at $h$ plus its total incoming
dimension.  The total dimension added at $h-1$ is at most $c_h$
times the largest ambient $d_h$ encountered.  Starting at the
highest level and descending gives polynomial bounds at each
level, because the depth is fixed.  These bounds are structural
and are fixed before choosing numerical accuracies.  In particular,
dimensions added at a given level use only higher-level dimensions.

The rational heights do not undergo one polynomial exponentiation
per cut.  At a fixed level the starting and incoming value spaces
are full coordinate spaces.  Stack all previously imposed integer
rows, extending old rows by zeros as copies arrive.  The current
$W_h$ is their joint kernel.  To choose the next section, solve this
stacked system together with $\lambda e=1$.  A nonzero minor and
Cramer's rule bound both the section and the rational kernel
height polynomially in the total dimension, cut count, and
$\log R_h$.  This is the simultaneous calculation in
Lemma~\ref{lifting:prepare-layer}, with the simpler full incoming
spaces here.  Lower polynomial substitutions and uniform scalar
residue rescalings do not change this value-space description.

We spell out the absence of a width circularity.  At a level $j$,
first fix all requirements inherited from levels above it.  Using
only those data and dimensions, choose $a_j$, the ratio $\upsilon_j/w_j$,
and the minimum-side threshold for omission.  The latter must
cover higher-level parent cutoffs and the separation scales in
Lemma~\ref{ap:major:backward}.  In that lemma the evaluated bounds
use only higher data, not the height or width of $K_j$.
Auxiliary cold threshold increases inside
Lemma~\ref{ap:major:refinement} merely remove further pre-short
directions temporarily; they do not change this warm minimum-side
choice.

Forecast a per-cut relative log-width loss $A_j$ using the same
data, without the initial $Q_j$.  A cut at $j$ uses only its own
dimension and gap factors.  A cut above $j$ may charge the upper
failure's coefficient, denominator, section, and width bounds,
the already fixed deviation ratios and spatial shrinkages, and
the omission cutoff.  These are warm at $j$.  If a common lower
width must lie below an upper width, multiply its current width
by the required small upper-scale factor; since all widths are
less than $1$, this also supplies the required absolute bound.
A cut below $j$ causes no shrinkage at $j$.  With $N_c$ the
structural cut bound, the preparation width is therefore at most
$Q_j+N_cA_j$ in logarithmic units.

Only now choose cold-pre and late budgets at $j$, permitting this
additive width bound.  They cover the bounded-depth major
calculations, denominator witness tests, perturbations, ranges,
rank thresholds, and parent cutoffs.  They also record the
requirements for lifting a failure at $j$ through its ancestors.
All calls at an ancestor can forecast the required pre ranges
before that ancestor's late choices.  Because there are only
fixedly many layer levels and bounded-depth calculations within
each call, this is a bounded downward composition of polynomial
log bounds, not a recursion through numerical precision at the
same level.

Choose gaps and individual target discounts using $N_c$ and the
dimension forecasts so that their total is as small as prescribed.
The positive selections are geometric coverage arguments and
therefore do not require a numerical lower bound on the initial
strict surplus.  Each spatial transaction has a polynomial log
cost in the full budgets.  Summing over at most $N_c$ such costs
and imposing the corresponding initial root-side cutoff sustains
all reindexings and subdivisions.  The finite cut bound then
forces termination with the stated probability tests and costs.
\end{proof}
